\documentclass[11pt,a4paper]{article}

% --- Codificación e idioma ---
\usepackage[utf8]{inputenc}
\usepackage[T1]{fontenc}
\usepackage[spanish,es-tabla,es-noshorthands]{babel}
\usepackage{lmodern}

% --- Geometría y formato general ---
\usepackage[margin=2.5cm]{geometry}
\usepackage{microtype}
\usepackage{parskip}
\usepackage{enumitem}
\setlist{topsep=2pt,itemsep=2pt,parsep=0pt}

% --- Tablas y diagramas ---
\usepackage{array}
\usepackage{booktabs}
\usepackage{amsmath,amssymb}
\usepackage{tikz}
\usetikzlibrary{positioning,arrows.meta,calc,shapes.geometric}

% --- Cabeceras y numeración ---
\usepackage{fancyhdr}
\pagestyle{fancy}
\fancyhf{}
\fancyhead[L]{\small Bases de Datos II}
\fancyhead[R]{\small Trabajo Práctico N.\textsuperscript{o} 7}
\fancyfoot[C]{\thepage}
\renewcommand{\headrulewidth}{0.4pt}

% --- Color y código ---
\usepackage{xcolor}
\definecolor{codebg}{RGB}{248,248,248}
\definecolor{codeframe}{RGB}{200,200,200}
\definecolor{codekw}{RGB}{0,0,180}
\definecolor{codestr}{RGB}{163,21,21}
\definecolor{codecmt}{RGB}{0,128,0}
\definecolor{notebg}{RGB}{255,250,230}
\definecolor{noteframe}{RGB}{220,180,80}

\usepackage{listings}

% Estilo para mongosh (JavaScript)
\lstdefinestyle{jsstyle}{
    language=Java,
    morekeywords={db,find,aggregate,insertOne,insertMany,updateOne,updateMany,deleteMany,countDocuments,createIndex,createCollection,dropDatabase,use,let,const,function},
    basicstyle=\ttfamily\footnotesize,
    keywordstyle=\color{codekw}\bfseries,
    stringstyle=\color{codestr},
    commentstyle=\color{codecmt}\itshape,
    backgroundcolor=\color{codebg},
    frame=single,
    rulecolor=\color{codeframe},
    breaklines=true,
    columns=fullflexible,
    keepspaces=true,
    showstringspaces=false,
    xleftmargin=0.4em,
    xrightmargin=0.4em,
    aboveskip=8pt,
    belowskip=8pt,
    tabsize=2,
    morecomment=[l]{//}
}

% Estilo para SQL (MySQL JSON)
\lstdefinelanguage{SQLext}[]{SQL}{
    morekeywords={JSON,JSON\_OBJECT,JSON\_ARRAY,JSON\_ARRAYAGG,JSON\_OBJECTAGG,
        JSON\_EXTRACT,JSON\_VALUE,JSON\_UNQUOTE,JSON\_CONTAINS,
        JSON\_SET,JSON\_TABLE,JSON\_TYPE,JSON\_LENGTH,SUBSTRING\_INDEX,
        VIRTUAL,STORED,GENERATED,ALWAYS,AS,PATH,COLUMNS,
        VARCHAR,DECIMAL,INT,DATE,DATETIME,TIMESTAMP,
        CREATE,TABLE,INSERT,INTO,VALUES,SELECT,FROM,WHERE,
        ORDER,BY,JOIN,ON,AND,OR,NOT,NULL,PRIMARY,KEY,
        FOREIGN,REFERENCES,INDEX,UNIQUE,WITH},
    sensitive=false,
    morecomment=[l]{--},
    morestring=[b]'
}
\lstdefinestyle{sqlstyle}{
    language=SQLext,
    basicstyle=\ttfamily\footnotesize,
    keywordstyle=\color{codekw}\bfseries,
    stringstyle=\color{codestr},
    commentstyle=\color{codecmt}\itshape,
    backgroundcolor=\color{codebg},
    frame=single,
    rulecolor=\color{codeframe},
    breaklines=true,
    columns=fullflexible,
    keepspaces=true,
    showstringspaces=false,
    xleftmargin=0.4em,
    xrightmargin=0.4em,
    aboveskip=8pt,
    belowskip=8pt,
    tabsize=2
}

% Estilo Java (Ej. 2 con Gson)
\lstdefinestyle{javastyle}{
    language=Java,
    basicstyle=\ttfamily\footnotesize,
    keywordstyle=\color{codekw}\bfseries,
    stringstyle=\color{codestr},
    commentstyle=\color{codecmt}\itshape,
    backgroundcolor=\color{codebg},
    frame=single,
    rulecolor=\color{codeframe},
    breaklines=true,
    columns=fullflexible,
    keepspaces=true,
    showstringspaces=false,
    xleftmargin=0.4em,
    xrightmargin=0.4em,
    aboveskip=8pt,
    belowskip=8pt,
    tabsize=2
}

% Estilo plano (JsonPath, JSON literal, shell)
\lstdefinestyle{plain}{
    basicstyle=\ttfamily\footnotesize,
    backgroundcolor=\color{codebg},
    frame=single,
    rulecolor=\color{codeframe},
    breaklines=true,
    columns=fullflexible,
    keepspaces=true,
    xleftmargin=0.4em,
    xrightmargin=0.4em,
    aboveskip=8pt,
    belowskip=8pt
}

\lstset{style=jsstyle}

% --- Cajas de nota ---
\usepackage[most]{tcolorbox}
\newtcolorbox{nota}{colback=notebg,colframe=noteframe,arc=2pt,
    left=6pt,right=6pt,top=4pt,bottom=4pt,boxrule=0.6pt}

% --- Hipervínculos ---
\usepackage[hidelinks,bookmarks=true]{hyperref}

% --- Comandos auxiliares (mismos que en Prácticos 1-6) ---
\newcommand{\itemcabecera}[1]{\noindent\textbf{#1}\par\medskip}
\newcommand{\bloqueConsigna}{\itemcabecera{Consigna}}
\newcommand{\bloqueIdea}{\itemcabecera{Idea}}
\newcommand{\bloqueSolucion}{\itemcabecera{Soluci\'on}}
\newcommand{\bloquePorque}{\itemcabecera{Por qu\'e funciona}}
\newcommand{\bloqueMySQL}{\itemcabecera{En MySQL 8+ (JSON)}}

\title{Trabajo Práctico N.\textsuperscript{o} 7: JSON y MongoDB\\[4pt]
       \large Bases de Datos II}
\author{Tomás Rodeghiero}
\date{2026}

\begin{document}

% =====================================================
% Carátula
% =====================================================
\begin{titlepage}
    \centering
    \vspace*{1.5cm}

    {\large \textbf{Universidad Nacional de Río Cuarto}}\\[6pt]
    {\normalsize Facultad de Ciencias Exactas, Físico-Químicas y Naturales}\\[4pt]
    {\normalsize Departamento de Computación}\\[4pt]
    {\normalsize Licenciatura en Ciencias de la Computación}\\[3cm]

    {\Large \textbf{Bases de Datos II}}\\[1.2cm]

    {\Large \textbf{Trabajo Práctico N.\textsuperscript{o} 7}}\\[6pt]
    {\large JSON y MongoDB}\\[3cm]

    \begin{tabular}{rl}
        \textbf{Alumno:} & Tomás Rodeghiero \\[4pt]
        \textbf{Año:}    & 2026 \\
    \end{tabular}

    \vfill
    {\small Resolución basada en el Teórico 9 (XML y JSON) y en experimentación con JsonPath/Gson en Java, MongoDB y Compass.}
\end{titlepage}

\tableofcontents
\newpage

% =====================================================
\section{Introducción}
% =====================================================

El Práctico 7 introduce dos tecnologías centrales del modelo documental: \textbf{JSON} como formato semiestructurado y \textbf{MongoDB} como motor de base de datos orientado a documentos. Los cinco ejercicios cubren ambos lados:

\begin{itemize}
    \item \textbf{Ejercicios 1 y 2}: trabajo con JSON desde Java. Consultas a un documento JSON con JsonPath (estilo XPath para JSON) y generación programática de un documento JSON a partir de una base relacional con la biblioteca Gson.
    \item \textbf{Ejercicios 3 a 5}: modelado en MongoDB. Diseño de colecciones, carga de datos y consultas (incluyendo \texttt{\$lookup} para joins entre colecciones).
\end{itemize}

Cada ejercicio sigue siempre la misma estructura: \textbf{Consigna} (texto literal del PDF), \textbf{Idea} (lectura natural y concepto detrás), \textbf{Solución} (JsonPath, Java, mongosh según corresponda), \textbf{Por qué funciona} (lógica paso a paso) y \textbf{En MySQL 8+ (JSON)} (equivalente con columnas \texttt{JSON} y funciones nativas: \texttt{JSON\_EXTRACT}, \texttt{JSON\_TABLE}, \texttt{JSON\_OBJECT}, \texttt{JSON\_ARRAYAGG}, operadores \texttt{->} y \texttt{->>}).

\begin{nota}
\textbf{Recursos del práctico.} La cátedra provee un proyecto Java con Gson y JsonPath, copiado en \texttt{practico-07/recursos/json/}. Los documentos de ejemplo son \texttt{json.txt} (de la teoría) y \texttt{jsonInventario.txt} (del Ejercicio~1). Para MongoDB se asume \texttt{mongosh} 1.x o superior conectando a un servidor 6.x+.
\end{nota}

% =====================================================
\section{Ejercicio 1: Consultas JsonPath sobre el inventario}
% =====================================================

\bloqueConsigna

Dado el documento \texttt{jsonInventario.txt} realice las siguientes consultas:

\begin{enumerate}[label=\alph*),leftmargin=*]
    \item Todos los items del inventario.
    \item Nombres de todos items del inventario.
    \item Item que están bajos en stock.
    \item Colores de los items cuyo alto es más de 10 cm.
\end{enumerate}

\bloqueIdea

JsonPath es un lenguaje de consulta para documentos JSON inspirado en XPath: \texttt{\$} representa la raíz, \texttt{.} navega hijos, \texttt{[*]} es comodín de arreglo y \texttt{[?(...)]} es un predicado de filtro. El concepto detrás del ejercicio es que un documento JSON anidado se puede explorar y filtrar declarativamente, igual que XPath sobre XML, sin necesidad de recorrer el árbol con código imperativo.

\bloqueSolucion

\textit{Estructura del documento:}

\begin{lstlisting}[style=plain]
{
  "inventario": {
    "items": [
      { "nombre": "tv",       "cantidad": 25,  "status": "A",
        "tamanio": { "alto": 14,    "ancho": 21,    "unidad": "cm" },
        "colores": ["negro","amarillo"] },
      { "nombre": "notebook", "cantidad": 50,  "status": "A",
        "tamanio": { "alto": 8.5,   "ancho": 11,    "unidad": "pulgadas" },
        "colores": ["rojo","negro"] },
      { "nombre": "mesa",     "cantidad": 100, "status": "D", ... },
      { "nombre": "planner",  "cantidad": 75,  "status": "D", ... },
      { "nombre": "postcard", "cantidad": 45,  "status": "A",
        "tamanio": { "alto": 10,    "ancho": 15.25, "unidad": "cm" },
        "colores": ["azul"] }
    ]
  },
  "stock_min": 60
}
\end{lstlisting}

Observaciones útiles: \texttt{stock\_min} está en la raíz (no dentro de cada ítem); la unidad varía entre \texttt{cm} y \texttt{pulgadas} (hay que filtrar por unidad para comparar alturas).

\textit{a) Todos los ítems:}

\begin{lstlisting}[style=plain]
$.inventario.items[*]
\end{lstlisting}

\textit{b) Nombres de todos los ítems:}

\begin{lstlisting}[style=plain]
$.inventario.items[*].nombre
-- Resultado: ["tv","notebook","mesa","planner","postcard"]
\end{lstlisting}

\textit{c) Ítems con bajo stock} ($\texttt{cantidad} < \texttt{stock\_min}$):

\begin{lstlisting}[style=plain]
$.inventario.items[?(@.cantidad < $['stock_min'])]
-- Cumplen: tv (25), notebook (50), postcard (45)
\end{lstlisting}

\textit{d) Colores de ítems con alto $>$ 10 cm} (filtrar por unidad y alto):

\begin{lstlisting}[style=plain]
$.inventario.items[?(@.tamanio.unidad == 'cm' && @.tamanio.alto > 10)].colores[*]
-- Solo tv (alto=14, cm) cumple; postcard tiene alto=10 (estricto, no entra)
-- Resultado: ["negro","amarillo"]
\end{lstlisting}

\textit{Ejecución desde Java con Jayway:}

\begin{lstlisting}[style=plain]
javac -cp "lib/*" -d bin src/jsonPath/Ejercicio1Inventario.java
java  -cp "bin:lib/*" jsonPath.Ejercicio1Inventario
\end{lstlisting}

\bloquePorque

\begin{itemize}
    \item En el inciso (a), \texttt{[*]} expande el arreglo \texttt{items} a todos sus elementos; sin el \texttt{[*]}, devolvería el array literal sin desempaquetar.
    \item En (b), la proyección \texttt{.nombre} aplicada a cada elemento del arreglo recolecta solo ese campo en un arreglo nuevo.
    \item En (c), el predicado \texttt{[?(...)]} compara la propiedad del nodo actual (\texttt{@.cantidad}) contra una referencia absoluta a la raíz (\texttt{\$['stock\_min']}). El estilo entre corchetes con comillas evita conflictos con el guión bajo o palabras reservadas.
    \item En (d), la conjunción de dos condiciones (\texttt{\&\&}) y la proyección final de \texttt{.colores[*]} achatan el array de arrays en un único arreglo de strings.
\end{itemize}

\bloqueMySQL

MySQL 8+ soporta el tipo \texttt{JSON} nativo, los operadores \texttt{->} (extracción) y \texttt{->>} (extracción + unquote), y \texttt{JSON\_TABLE} para ``aplanar'' arrays a filas relacionales.

\begin{lstlisting}[style=sqlstyle]
-- Tabla con el documento como columna JSON
CREATE TABLE inventario_doc (
    id INT PRIMARY KEY,
    datos JSON NOT NULL
);

INSERT INTO inventario_doc (id, datos)
VALUES (1, '{"inventario":{"items":[...]},"stock_min":60}');

-- a) Todos los items del inventario
SELECT datos->>'$.inventario.items' AS items FROM inventario_doc;

-- b) Nombres de todos los items (JSON_TABLE aplana el array)
SELECT it.nombre
FROM inventario_doc d,
     JSON_TABLE(
         d.datos, '$.inventario.items[*]'
         COLUMNS (nombre VARCHAR(100) PATH '$.nombre')
     ) AS it;

-- c) Items bajos en stock (cantidad < stock_min)
SELECT it.nombre, it.cantidad
FROM inventario_doc d,
     JSON_TABLE(
         d.datos, '$.inventario.items[*]'
         COLUMNS (
             nombre   VARCHAR(100) PATH '$.nombre',
             cantidad INT          PATH '$.cantidad'
         )
     ) AS it
WHERE it.cantidad < JSON_VALUE(d.datos, '$.stock_min');

-- d) Colores de items con alto > 10 cm
SELECT it.nombre, it.colores
FROM inventario_doc d,
     JSON_TABLE(
         d.datos, '$.inventario.items[*]'
         COLUMNS (
             nombre  VARCHAR(100)   PATH '$.nombre',
             alto    DECIMAL(10,2)  PATH '$.tamanio.alto',
             unidad  VARCHAR(20)    PATH '$.tamanio.unidad',
             colores JSON           PATH '$.colores'
         )
     ) AS it
WHERE it.unidad = 'cm' AND it.alto > 10;
\end{lstlisting}

\textit{Diferencias relevantes con JsonPath puro:}

\begin{itemize}
    \item \texttt{JSON\_TABLE} es la herramienta clave: convierte un array JSON en filas SQL para poder usar \texttt{WHERE}, \texttt{ORDER BY} y joins habituales. Sin esto, hay que combinar varios \texttt{JSON\_EXTRACT} con índices duros.
    \item Los operadores \texttt{->} y \texttt{->>} solo funcionan con \emph{paths fijos} (no admiten predicados como \texttt{[?(...)]}); por eso para condiciones se necesita \texttt{JSON\_TABLE}.
    \item Para indexar accesos JSON frecuentes en MySQL hay que crear \emph{columnas generadas} (\texttt{GENERATED ALWAYS AS (...) STORED}) e indexar esas columnas; los índices funcionales sobre \texttt{JSON\_EXTRACT} llegaron recién en 8.0.13.
\end{itemize}

% =====================================================
\section{Ejercicio 2: Generar un JSON único de facturas e items}
% =====================================================

\bloqueConsigna

Realice un programa Java (ver biblioteca GSON) que cree un único documento json con la información almacenada en la tabla Factura y sus Items (base de datos del ejercicio 1) a de la práctica 1).

\bloqueIdea

Se pide \emph{aplanar} un modelo relacional (1:N entre \texttt{factura} e \texttt{item\_factura}) a un único documento JSON con un arreglo de facturas y, dentro de cada factura, un arreglo anidado de items. Es el patrón canónico de exportación a un modelo documental; Gson se ocupa de la serialización segura (escape de comillas, codificación de números/fechas, manejo de \texttt{null}).

\bloqueSolucion

\textit{Esquema del documento generado:}

\begin{lstlisting}[style=plain]
{
  "generado_en": "2026-03-28T12:34:56.123-03:00",
  "facturas": [
    {
      "nro_factura": 1001,
      "nro_cliente": 1,
      "fecha": "2024-08-01",
      "monto": 106000.0,
      "items": [
        { "cod_producto": 102, "cantidad": 2, "precio":  8000.0 },
        { "cod_producto": 103, "cantidad": 1, "precio": 90000.0 }
      ]
    },
    ...
  ]
}
\end{lstlisting}

\textit{Estructura del proyecto Java:}

\begin{itemize}
    \item \texttt{recursos/json/src/jsonPath/Ejercicio2FacturaJson.java}: main del programa.
    \item \texttt{recursos/json/config/db.mysql.properties}: parámetros de conexión y nombres de tablas.
    \item \texttt{recursos/json/lib/gson-2.8.6.jar}, \texttt{mysql-connector-j-8.0.33.jar}: dependencias.
\end{itemize}

\textit{Flujo del programa:}

\begin{enumerate}
    \item lee la configuración del \texttt{.properties};
    \item abre una conexión JDBC a \texttt{practico1a\_mysql};
    \item lee todas las facturas en un \texttt{Map<Long, Factura>};
    \item lee los items y los agrupa por \texttt{nro\_factura};
    \item arma una raíz \texttt{\{ generado\_en, facturas[] \}};
    \item serializa con \texttt{new GsonBuilder().setPrettyPrinting().create().toJson(raiz)};
    \item escribe el resultado en \texttt{salida/facturas\_items.json}.
\end{enumerate}

\textit{Bloque clave del programa (Gson):}

\begin{lstlisting}[style=javastyle]
import com.google.gson.Gson;
import com.google.gson.GsonBuilder;

Gson gson = new GsonBuilder()
    .setPrettyPrinting()
    .setDateFormat("yyyy-MM-dd")
    .create();

String json = gson.toJson(raiz);
Files.writeString(Path.of("salida/facturas_items.json"), json);
\end{lstlisting}

\textit{Compilación y ejecución:}

\begin{lstlisting}[style=plain]
javac -cp "lib/*" -d bin src/jsonPath/Ejercicio2FacturaJson.java
java  -cp "bin:lib/*" jsonPath.Ejercicio2FacturaJson config/db.mysql.properties
\end{lstlisting}

\bloquePorque

\begin{itemize}
    \item El paso 4 (agrupar items por \texttt{nro\_factura}) es donde ocurre el ``aplanado'': los items dejan de ser una tabla aparte y se anidan dentro de su factura.
    \item Gson maneja automáticamente el escape de strings (\texttt{"} dentro de un campo se convierte en \texttt{\textbackslash"}), las fechas y los \texttt{null}; construir el JSON con \texttt{StringBuilder} es propenso a errores en cuanto los datos tienen caracteres especiales.
    \item \texttt{setPrettyPrinting()} indenta el JSON con dos espacios para que sea legible; en producción se suele omitir para ahorrar espacio.
\end{itemize}

\bloqueMySQL

MySQL 8+ permite generar el mismo documento JSON \emph{directamente en SQL}, sin pasar por Java. \texttt{JSON\_OBJECT} arma un objeto, \texttt{JSON\_ARRAYAGG} agrupa filas en un arreglo y se pueden anidar para construir documentos arbitrarios:

\begin{lstlisting}[style=sqlstyle]
USE practico1a_mysql;

SELECT JSON_OBJECT(
    'generado_en', NOW(),
    'facturas', (
        SELECT JSON_ARRAYAGG(
            JSON_OBJECT(
                'nro_factura', f.nro_factura,
                'nro_cliente', f.nro_cliente,
                'fecha',       f.fecha,
                'monto',       f.monto,
                'items', (
                    SELECT JSON_ARRAYAGG(
                        JSON_OBJECT(
                            'cod_producto', i.cod_producto,
                            'cantidad',     i.cantidad,
                            'precio',       i.precio
                        )
                    )
                    FROM item_factura i
                    WHERE i.nro_factura = f.nro_factura
                )
            )
        )
        FROM factura f
    )
) AS documento;
\end{lstlisting}

\textit{Diferencias relevantes con el enfoque Java+Gson:}

\begin{itemize}
    \item Una sola sentencia SQL reemplaza el programa Java entero: se ahorra todo el código de conexión, agrupado en memoria y serialización.
    \item \texttt{JSON\_ARRAYAGG} no garantiza orden a menos que se use con una subconsulta \texttt{ORDER BY} (en MySQL 8.0.14+ se puede usar \texttt{... ORDER BY} dentro de la función agregada).
    \item Para volcar el resultado a archivo se usa \texttt{SELECT ... INTO OUTFILE} o se redirige la salida del cliente (\texttt{mysql -e "SELECT ..."  > facturas\_items.json}).
\end{itemize}

% =====================================================
\section{Ejercicio 3: Diseño de la base en MongoDB}
% =====================================================

\bloqueConsigna

Cree una base de datos en MongoDB que contengan colecciones que permitan cargar datos con la siguiente estructura definida con un diagrama de E/R:

\begin{itemize}
    \item \texttt{usuarios}: \texttt{nick}, \texttt{NyApe}, \texttt{edad}, \texttt{direccion} (compuesta por \texttt{calle} y \texttt{altura}), \texttt{telefonos}.
    \item \texttt{noticias}: \texttt{codigo}, \texttt{texto}, \texttt{fecha}; cada usuario \emph{crea} 1:N noticias.
    \item \texttt{comentarios}: \texttt{codigo}, \texttt{texto}; cada usuario \emph{hace} 1:N comentarios, y cada noticia tiene N comentarios \emph{para} ella, con palabras clave.
\end{itemize}

\textbf{Nota:} El campo dirección está compuesto por calle y altura.

\bloqueIdea

MongoDB es documental: el dato unitario es un documento JSON/BSON y se agrupa en \emph{colecciones} (equivalente lógico de tablas pero sin esquema rígido). La diferencia clave con el modelo relacional es la decisión \emph{embedding vs referencing}: atributos compuestos (\texttt{direccion}) y multivaluados (\texttt{telefonos}) van \textbf{embebidos} dentro del usuario; entidades con ciclo de vida independiente (\texttt{noticias}, \texttt{comentarios}) van en colecciones \textbf{separadas} con referencias por clave de negocio.

\bloqueSolucion

\textit{Decisión de modelado:}

\begin{itemize}
    \item \texttt{direccion} y \texttt{telefonos} van dentro de \texttt{usuarios} (atributos del usuario, no entidades).
    \item \texttt{noticias} y \texttt{comentarios} se mantienen como colecciones aparte porque su volumen y ciclo de vida son independientes.
\end{itemize}

\textit{Esquema de las colecciones:}

\begin{lstlisting}[style=plain]
// usuarios
{
  nick:       "<clave de negocio, unica>",
  nyape:      "<nombre y apellido>",
  edad:       <int>,
  direccion:  { calle: "<string>", altura: <int> },
  telefonos:  [ "<tel1>", "<tel2>", ... ]
}

// noticias
{
  codigo:      <int, unico>,
  texto:       "<string>",
  fecha:       <Date>,
  autor_nick:  "<ref logica a usuarios.nick>"
}

// comentarios
{
  codigo:          <int, unico>,
  texto:           "<string>",
  autor_nick:      "<ref logica a usuarios.nick>",
  noticia_codigo:  <int, ref logica a noticias.codigo>,
  pal_claves:      [ "<kw1>", "<kw2>", ... ]
}
\end{lstlisting}

\textit{Script de creación (mongosh) con \texttt{\$jsonSchema} e índices:}

\begin{lstlisting}
use("practico7_er_noticias");
db.dropDatabase();

db.createCollection("usuarios", {
  validator: { $jsonSchema: {
    bsonType: "object",
    required: ["nick","nyape","edad","direccion","telefonos"],
    properties: {
      nick:      { bsonType: "string", minLength: 1 },
      nyape:     { bsonType: "string", minLength: 1 },
      edad:      { bsonType: "int",    minimum:   0 },
      direccion: { bsonType: "object",
                   required: ["calle","altura"],
                   properties: {
                     calle:  { bsonType: "string", minLength: 1 },
                     altura: { bsonType: "int",    minimum:   0 }
                   }},
      telefonos: { bsonType: "array",  minItems: 1,
                   items: { bsonType: "string", minLength: 1 }}
    }
  }}
});

// ... similar para noticias y comentarios

db.usuarios.createIndex   ({ nick:           1 }, { unique: true });
db.noticias.createIndex   ({ codigo:         1 }, { unique: true });
db.comentarios.createIndex({ codigo:         1 }, { unique: true });
db.noticias.createIndex   ({ autor_nick:     1 });
db.comentarios.createIndex({ autor_nick:     1 });
db.comentarios.createIndex({ noticia_codigo: 1 });
\end{lstlisting}

\bloquePorque

\begin{itemize}
    \item \texttt{direccion} y \texttt{telefonos} embebidos: la consulta dominante es ``traer un usuario completo'', así se evita un join lógico y la cardinalidad por usuario es chica (1 dirección, pocos teléfonos).
    \item \texttt{comentarios} en colección aparte (en vez de embebidos en la noticia): los comentarios crecen sin tope con la actividad, mantenerlos embebidos haría documentos cada vez más grandes; además tienen \texttt{codigo} propio y \texttt{pal\_claves}, lo que los hace entidades de pleno derecho.
    \item \texttt{\$jsonSchema} aporta validación de tipos a la inserción, simulando los \texttt{CHECK} del modelo relacional sin perder la flexibilidad documental.
    \item Los índices únicos sobre \texttt{nick} y \texttt{codigo} aseguran la integridad de claves de negocio (MongoDB no tiene FKs, así que esta validación queda del lado de la aplicación).
\end{itemize}

\bloqueMySQL

El mismo E/R se puede modelar en MySQL 8+ con columnas \texttt{JSON}. Para indexar atributos dentro del JSON (\texttt{nick}, \texttt{autor\_nick}, \texttt{apellido}), se usan \emph{columnas generadas} \texttt{STORED} sobre las cuales sí se puede crear un \texttt{INDEX} o \texttt{FOREIGN KEY}:

\begin{lstlisting}[style=sqlstyle]
CREATE DATABASE practico7_er_mysql;
USE practico7_er_mysql;

CREATE TABLE usuarios (
    nick     VARCHAR(60) PRIMARY KEY,
    datos    JSON NOT NULL,
    apellido VARCHAR(60) GENERATED ALWAYS AS (
        SUBSTRING_INDEX(JSON_UNQUOTE(datos->>'$.nyape'), ' ', -1)
    ) STORED,
    edad     INT GENERATED ALWAYS AS (
        CAST(JSON_EXTRACT(datos, '$.edad') AS UNSIGNED)
    ) STORED,
    INDEX idx_apellido (apellido),
    INDEX idx_edad     (edad)
);

CREATE TABLE noticias (
    codigo      INT PRIMARY KEY,
    datos       JSON NOT NULL,
    autor_nick  VARCHAR(60) GENERATED ALWAYS AS (
        JSON_UNQUOTE(datos->>'$.autor_nick')
    ) STORED,
    fecha       DATE GENERATED ALWAYS AS (
        CAST(JSON_UNQUOTE(datos->>'$.fecha') AS DATE)
    ) STORED,
    INDEX idx_autor (autor_nick),
    INDEX idx_fecha (fecha),
    FOREIGN KEY (autor_nick) REFERENCES usuarios(nick)
);

CREATE TABLE comentarios (
    codigo          INT PRIMARY KEY,
    datos           JSON NOT NULL,
    autor_nick      VARCHAR(60) GENERATED ALWAYS AS (
        JSON_UNQUOTE(datos->>'$.autor_nick')
    ) STORED,
    noticia_codigo  INT GENERATED ALWAYS AS (
        CAST(JSON_EXTRACT(datos, '$.noticia_codigo') AS UNSIGNED)
    ) STORED,
    INDEX idx_autor   (autor_nick),
    INDEX idx_noticia (noticia_codigo),
    FOREIGN KEY (autor_nick)     REFERENCES usuarios(nick),
    FOREIGN KEY (noticia_codigo) REFERENCES noticias(codigo)
);
\end{lstlisting}

\textit{Limitaciones a notar:}

\begin{itemize}
    \item MySQL no tiene \emph{pipeline} de agregación como MongoDB: las consultas analíticas se escriben con SQL estándar más funciones JSON.
    \item No hay \texttt{\$lookup} recursivo (\texttt{\$graphLookup} en MongoDB): los joins son los de SQL clásico.
    \item Los índices sobre JSON requieren columnas generadas; un \texttt{INDEX} directo sobre \texttt{JSON\_EXTRACT} no se permite hasta crear primero la generated column.
\end{itemize}

% =====================================================
\section{Ejercicio 4: Carga mínima}
% =====================================================

\bloqueConsigna

Agregue a la base de datos creada al menos dos usuarios, 4 noticias y dos comentarios por noticia.

\bloqueIdea

Es una carga seed. La sutileza está en hacerla \emph{idempotente}: los índices únicos creados en el Ejercicio 3 fallarían si se re-ejecutan los \texttt{insertOne}. El idiom de MongoDB es \texttt{updateOne(..., \{ upsert: true \})}: si el documento existe se actualiza, si no se inserta.

\bloqueSolucion

\textit{Datos a insertar:}

\begin{itemize}
    \item 2 usuarios: \texttt{moises}, \texttt{elias}.
    \item 4 noticias con \texttt{codigo}\,$\in\{2001,2002,2003,2004\}$.
    \item 8 comentarios (2 por noticia) con \texttt{codigo}\,$\in\{7001,\ldots,7008\}$.
\end{itemize}

\textit{Patrón \texttt{upsert} para idempotencia:}

\begin{lstlisting}
db.usuarios.updateOne(
  { nick: "moises" },
  { $set: {
      nyape: "Moises Gonzalez",
      edad:  NumberInt(45),
      direccion: { calle: "San Martin", altura: NumberInt(123) },
      telefonos: ["358-1111111","358-2222222"]
  }},
  { upsert: true }
);

db.usuarios.updateOne(
  { nick: "elias" },
  { $set: {
      nyape: "Elias Perez",
      edad:  NumberInt(38),
      direccion: { calle: "Belgrano", altura: NumberInt(450) },
      telefonos: ["358-3333333"]
  }},
  { upsert: true }
);

// 4 noticias (codigo 2001..2004), autor_nick alternando moises/elias
db.noticias.updateOne(
  { codigo: 2001 },
  { $set: {
      texto: "Nueva politica de privacidad",
      fecha: ISODate("2024-08-01"),
      autor_nick: "moises"
  }},
  { upsert: true }
);
// ... 2002, 2003, 2004 con autor_nick variable

// 8 comentarios (2 por noticia)
db.comentarios.updateOne(
  { codigo: 7001 },
  { $set: {
      texto: "Muy buena noticia",
      autor_nick: "elias",
      noticia_codigo: 2001,
      pal_claves: ["positivo","actualidad"]
  }},
  { upsert: true }
);
// ... 7002..7008
\end{lstlisting}

\textit{Verificación de conteos:}

\begin{lstlisting}
db.usuarios.countDocuments({ nick: { $in: ["moises","elias"] } });       // 2
db.noticias.countDocuments({ codigo: { $in: [2001,2002,2003,2004] } });  // 4
db.comentarios.countDocuments({ noticia_codigo: { $in: [2001,2002,2003,2004] } }); // 8

// Confirmar 2 comentarios por noticia
db.comentarios.aggregate([
  { $match: { noticia_codigo: { $in: [2001,2002,2003,2004] } } },
  { $group: { _id: "$noticia_codigo", total: { $sum: 1 } } },
  { $sort:  { _id: 1 } }
]);
// 4 filas, todas con total=2
\end{lstlisting}

\bloquePorque

\begin{itemize}
    \item \texttt{updateOne} con \texttt{upsert: true} es el patrón estándar para cargas idempotentes en MongoDB: si la fila clave (\texttt{nick} para usuarios, \texttt{codigo} para noticias y comentarios) ya existe, se actualiza; si no, se inserta. Re-correr el script no rompe nada.
    \item \texttt{NumberInt(...)} fuerza el tipo \texttt{int} de BSON; sin esto, mongosh interpreta los números como \texttt{double}, lo que choca con el validador \texttt{\$jsonSchema} que pide \texttt{bsonType: "int"}.
    \item \texttt{ISODate("...")} crea un valor \texttt{Date} de BSON; si se cargan strings la consulta del Ejercicio 5.c con \texttt{\$sort} podría ordenar lexicográficamente en vez de cronológicamente.
\end{itemize}

\bloqueMySQL

El equivalente MySQL usa \texttt{INSERT ... ON DUPLICATE KEY UPDATE} para mantener idempotencia. Los documentos JSON se escriben como literales de string:

\begin{lstlisting}[style=sqlstyle]
USE practico7_er_mysql;

INSERT INTO usuarios (nick, datos) VALUES
  ('moises', JSON_OBJECT(
      'nyape',     'Moises Gonzalez',
      'edad',      45,
      'direccion', JSON_OBJECT('calle','San Martin','altura',123),
      'telefonos', JSON_ARRAY('358-1111111','358-2222222')
  )),
  ('elias',  JSON_OBJECT(
      'nyape',     'Elias Perez',
      'edad',      38,
      'direccion', JSON_OBJECT('calle','Belgrano','altura',450),
      'telefonos', JSON_ARRAY('358-3333333')
  ))
ON DUPLICATE KEY UPDATE datos = VALUES(datos);

INSERT INTO noticias (codigo, datos) VALUES
  (2001, JSON_OBJECT('texto','Nueva politica de privacidad',
                     'fecha','2024-08-01','autor_nick','moises')),
  (2002, JSON_OBJECT('texto','Resultados del torneo',
                     'fecha','2024-08-05','autor_nick','elias')),
  (2003, JSON_OBJECT('texto','Apertura del museo',
                     'fecha','2024-08-10','autor_nick','moises')),
  (2004, JSON_OBJECT('texto','Cambio en el transporte',
                     'fecha','2024-08-15','autor_nick','elias'))
ON DUPLICATE KEY UPDATE datos = VALUES(datos);

-- comentarios (8 filas, 2 por noticia)
INSERT INTO comentarios (codigo, datos) VALUES
  (7001, JSON_OBJECT('texto','Muy buena noticia','autor_nick','elias',
                     'noticia_codigo',2001,'pal_claves',JSON_ARRAY('positivo','actualidad'))),
  -- ... 7002..7008
ON DUPLICATE KEY UPDATE datos = VALUES(datos);
\end{lstlisting}

\textit{Diferencia clave con MongoDB:} las columnas generadas (\texttt{apellido}, \texttt{autor\_nick}, etc.) se calculan automáticamente con cada \texttt{INSERT}/\texttt{UPDATE}, así que no hay que duplicar la lógica de extracción al cargar datos. \texttt{ON DUPLICATE KEY UPDATE} reemplaza el documento completo cuando la PK ya existe (equivalente a \texttt{upsert}).

% =====================================================
\section{Ejercicio 5: Consultas MongoDB}
% =====================================================

\bloqueConsigna

Resuelva las siguientes consultas:

\begin{enumerate}[label=\alph*),leftmargin=*]
    \item Listar el nick, nombre y apellido y edad de todos los usuarios.
    \item Listar los usuarios de más de 40 años.
    \item Devolver las noticias que crearon los usuarios de apellido Gonzalez ordenadas por fecha.
\end{enumerate}

\bloqueIdea

Tres consultas con dificultad creciente. La (a) requiere transformar un campo concatenado (\texttt{nyape}) en dos campos separados; el aggregation pipeline lo hace con \texttt{\$split} y \texttt{\$arrayElemAt}. La (b) es un \texttt{find} con filtro y proyección. La (c) requiere un \emph{join lógico} entre \texttt{noticias} y \texttt{usuarios} (porque \texttt{autor\_nick} es referencia, no embedding), que en MongoDB se resuelve con \texttt{\$lookup} + \texttt{\$unwind}.

\bloqueSolucion

\textit{a) Listar nick, nombre, apellido, edad:}

\begin{lstlisting}
db.usuarios.aggregate([
  {
    $project: {
      _id: 0,
      nick: 1,
      edad: 1,
      nombre:   { $arrayElemAt: [{ $split: ["$nyape", " "] },  0] },
      apellido: { $arrayElemAt: [{ $split: ["$nyape", " "] }, -1] }
    }
  }
]);
\end{lstlisting}

\texttt{\$split} parte \texttt{nyape} por espacios; \texttt{\$arrayElemAt} con índice 0 toma el primer token (nombre) y con -1 toma el último (apellido). La convención cubre nombres compuestos como "Juan Carlos Perez".

\textit{b) Usuarios de más de 40 años:}

\begin{lstlisting}
db.usuarios.find(
  { edad: { $gt: 40 } },
  { _id: 0, nick: 1, nyape: 1, edad: 1 }
);
\end{lstlisting}

\textit{c) Noticias de autores apellido Gonzalez ordenadas por fecha:}

\begin{lstlisting}
db.noticias.aggregate([
  {
    $lookup: {
      from:         "usuarios",
      localField:   "autor_nick",
      foreignField: "nick",
      as:           "autor"
    }
  },
  { $unwind: "$autor" },                       // $lookup devuelve array
  { $match: { "autor.nyape": /gonzalez$/i } }, // apellido termina en Gonzalez
  { $sort:  { fecha: 1 } },
  {
    $project: {
      _id: 0,
      codigo: 1,
      texto: 1,
      fecha: 1,
      autor_nick: 1,
      autor: "$autor.nyape"
    }
  }
]);
\end{lstlisting}

\bloquePorque

\begin{itemize}
    \item En (a), \texttt{\$split} + \texttt{\$arrayElemAt} reproduce la lógica de \texttt{SUBSTRING\_INDEX} de MySQL pero dentro del pipeline. El \texttt{-1} en \texttt{\$arrayElemAt} es el truco para tomar el último elemento sin conocer el largo.
    \item En (b) basta con \texttt{find}: el filtro \texttt{\{ edad: \{ \$gt: 40 \} \}} es directo y la proyección suprime \texttt{\_id} (que MongoDB siempre incluye por defecto).
    \item En (c) el pipeline tiene cinco etapas: (1) \texttt{\$lookup} simula el join por igualdad; (2) \texttt{\$unwind} convierte el array de un solo elemento en un objeto plano; (3) \texttt{\$match} filtra con regex insensible a mayúsculas; (4) \texttt{\$sort} ordena cronológicamente (requiere que \texttt{fecha} sea \texttt{Date}, no string); (5) \texttt{\$project} arma el resultado.
    \item Para que el \texttt{\$lookup} sea eficiente con grandes volúmenes, conviene tener \texttt{usuarios.nick} indexado (lo cumple el índice único del Ejercicio 3).
\end{itemize}

\bloqueMySQL

Las tres consultas se traducen a SQL con funciones JSON. Los índices sobre columnas generadas (\texttt{apellido}, \texttt{autor\_nick}, \texttt{fecha}) hacen que los filtros y \texttt{ORDER BY} sean eficientes:

\begin{lstlisting}[style=sqlstyle]
-- a) nick, nombre, apellido, edad
SELECT
    nick,
    SUBSTRING_INDEX(JSON_UNQUOTE(datos->>'$.nyape'), ' ',  1) AS nombre,
    SUBSTRING_INDEX(JSON_UNQUOTE(datos->>'$.nyape'), ' ', -1) AS apellido,
    JSON_EXTRACT(datos, '$.edad')                             AS edad
FROM usuarios;

-- b) Usuarios de mas de 40 anios (usa el indice idx_edad sobre la columna generada)
SELECT
    nick,
    datos->>'$.nyape' AS nyape,
    edad
FROM usuarios
WHERE edad > 40;

-- c) Noticias de autores apellido Gonzalez ordenadas por fecha
SELECT
    n.codigo,
    n.datos->>'$.texto'      AS texto,
    n.fecha,
    n.autor_nick,
    u.datos->>'$.nyape'      AS autor
FROM noticias n
JOIN usuarios u ON u.nick = n.autor_nick
WHERE u.apellido = 'Gonzalez'
ORDER BY n.fecha;
\end{lstlisting}

\textit{Diferencias relevantes con MongoDB:}

\begin{itemize}
    \item \texttt{\$lookup + \$unwind} se reemplaza por un \texttt{JOIN} clásico. Es más simple y, con índices, más eficiente: el optimizador puede usar \emph{hash join} o \emph{merge join}.
    \item \texttt{\$split} + \texttt{\$arrayElemAt} se reemplaza por \texttt{SUBSTRING\_INDEX(..., ' ', -1)}. Más conciso, pero menos expresivo si el formato del campo concatenado fuera más complejo.
    \item Para que la regex \texttt{/gonzalez\$/i} de MongoDB sea equivalente, en MySQL se puede usar \texttt{REGEXP '(?i)gonzalez\$'} en lugar de la igualdad estricta (\texttt{apellido = 'Gonzalez'}); la igualdad asume que la columna generada normalizó la capitalización.
    \item MySQL \emph{no} tiene un equivalente directo del aggregation pipeline; los reportes analíticos se hacen con SQL + \texttt{WITH} (CTEs) o con \texttt{GROUP BY} clásico, y son menos componibles que la cadena de etapas de MongoDB.
\end{itemize}

% =====================================================
\section*{Conclusiones}
\addcontentsline{toc}{section}{Conclusiones}
% =====================================================

El práctico recorre tres capas del trabajo con datos semiestructurados:

\begin{enumerate}
    \item \textbf{Consultas sobre JSON} (Ejercicio 1): JsonPath es a JSON lo que XPath es a XML. Ofrece operadores de navegación (\texttt{.}, \texttt{[*]}, \texttt{..}), filtrado (\texttt{[?(...)]}) y agregaciones.
    \item \textbf{Generación de JSON desde código} (Ejercicio 2): el caso típico es leer un modelo relacional y emitir un documento JSON \emph{aplanado}. Gson resuelve la serialización; la lógica de agrupado queda del lado de la aplicación. MySQL 8+ permite el mismo resultado en una sola sentencia SQL con \texttt{JSON\_OBJECT} y \texttt{JSON\_ARRAYAGG}.
    \item \textbf{Modelado documental con MongoDB} (Ejercicios 3--5): las decisiones de embedding vs referencing, las claves de negocio y los índices únicos son las herramientas para representar un E/R sin perder integridad. El aggregation framework cubre la mayoría de las consultas no triviales, incluyendo joins lógicos con \texttt{\$lookup}.
\end{enumerate}

La idea común detrás del práctico es: \emph{el modelo documental no es ``relacional sin reglas''}. Tiene sus propias buenas prácticas (embedding por composición, referencing por relaciones independientes, validación con \texttt{\$jsonSchema}, índices) y sus propias herramientas de consulta. Pero también es importante reconocer que \textbf{MySQL 8+ implementa JSON nativo} con funciones potentes: para muchos casos de uso donde no se necesitan replicación distribuida ni esquemas verdaderamente dinámicos, una tabla con columna \texttt{JSON} y columnas generadas indexadas es una alternativa razonable a MongoDB, conservando transacciones ACID y SQL relacional.

\end{document}
